\chapter{ROS 2 para principiantes}

\begin{enpocas}
Qué es ROS~2, la terminal en cinco minutos, nodos, tópicos, mensajes, servicios, parámetros,
archivos \emph{launch} y el espacio de trabajo. Con ejemplos de nuestra simulación.
\end{enpocas}

\section{La terminal en cinco minutos}

Casi todo se hace escribiendo órdenes en una \textbf{terminal} (Ctrl+Alt+T en Ubuntu). Lo mínimo:

\begin{lstlisting}[style=bash]
pwd                      # ¿en qué carpeta estoy?
ls                       # ¿qué hay aquí?
cd ~/ros2_ws             # entrar a una carpeta (~ es tu carpeta personal)
cd ..                    # subir un nivel
cat archivo.txt          # ver un archivo
nano archivo.txt         # editar un archivo (Ctrl+O guarda, Ctrl+X sale)
Ctrl+C                   # detener el programa que está corriendo
\end{lstlisting}

\begin{consejo}
La tecla Tab completa nombres de archivos y órdenes; la flecha hacia arriba recupera la orden
anterior. Ahorran muchísimo tiempo y errores de dedo.
\end{consejo}

\section{¿Qué es ROS 2?}

\textbf{ROS} (\emph{Robot Operating System}) no es un sistema operativo: es un conjunto de
herramientas para que muchos programas pequeños de un robot hablen entre sí. En lugar de un programa
gigante que lo hace todo, hay un programa que lee la cámara, otro que detecta fuego, otro que decide
a dónde ir, otro que controla los motores... y ROS los conecta. Usamos la versión \textbf{ROS~2
Jazzy}.

\begin{ejemplo}
Piensa en una cocina de restaurante. El que recibe pedidos los pega en un tablero; el cocinero toma
los pedidos del tablero; el mesero recoge los platos de la barra. Nadie habla directamente con nadie:
todos leen y escriben en lugares comunes. En ROS, los cocineros son los \emph{nodos} y los tableros
son los \emph{tópicos}.
\end{ejemplo}

\section{Nodos, tópicos y mensajes}

\begin{description}
  \item[Nodo:] un programa que hace una tarea. En nuestra simulación: \cod{detector_yolo},
        \cod{busqueda_incendio}, \cod{controlador_sta}, \cod{viento}...
  \item[Tópico:] un canal con nombre por donde viajan datos, por ejemplo \cod{/cuadrirotor/camara}.
  \item[Mensaje:] el dato que viaja, con un tipo definido. \cod{sensor_msgs/Image} es una imagen;
        \cod{nav_msgs/Odometry}, una posición con velocidad.
  \item[Publicar / suscribirse:] un nodo \emph{publica} en un tópico y cualquier otro se
        \emph{suscribe} para recibir. El que publica no sabe quién escucha.
\end{description}

\begin{figure}[H]
\centering
\begin{tikzpicture}[
  nodo/.style={draw,ellipse,fill=green!12,minimum width=2.8cm,minimum height=0.9cm,font=\small},
  top/.style={draw,rectangle,fill=yellow!20,minimum width=3.2cm,minimum height=0.7cm,font=\scriptsize\ttfamily},
  fl/.style={-{Stealth[length=2.2mm]},thick}]
  \node[nodo] (gz) at (0,0) {Gazebo + puente};
  \node[top] (cam) at (4.2,0.8) {/cuadrirotor/camara};
  \node[top] (odo) at (4.2,-0.8) {/cuadrirotor/odometria};
  \node[nodo] (det) at (8.6,0.8) {detector\_yolo};
  \node[top] (detc) at (8.6,-1.0) {/cuadrirotor/detecciones};
  \node[nodo] (bus) at (8.6,-2.8) {busqueda\_incendio};
  \node[top] (con) at (4.2,-2.8) {/cuadrirotor/consigna};
  \node[nodo] (ctl) at (0,-2.8) {controlador\_sta};
  \draw[fl] (gz) -- (cam); \draw[fl] (gz) -- (odo); \draw[fl] (cam) -- (det);
  \draw[fl] (det) -- (detc); \draw[fl] (detc) -- (bus); \draw[fl] (odo) -- (bus);
  \draw[fl] (bus) -- (con); \draw[fl] (con) -- (ctl); \draw[fl] (odo) -- (ctl);
  \draw[fl] (ctl) -- node[left,font=\scriptsize]{/cuadrirotor/motores} (gz);
\end{tikzpicture}
\caption{Nodos (óvalos verdes) que se comunican por tópicos (rectángulos amarillos) en la misión de
búsqueda.}
\end{figure}

Para mirar desde la terminal:
\begin{lstlisting}[style=bash]
ros2 node list                                   # nodos que están corriendo
ros2 topic list                                  # tópicos que existen
ros2 topic echo /cuadrirotor/busqueda/estado     # ver lo que pasa por un tópico
ros2 topic hz /cuadrirotor/camara                # ¿cuántos mensajes por segundo?
ros2 topic info /cuadrirotor/odometria           # tipo y quién publica/escucha
ros2 interface show nav_msgs/msg/Odometry        # ¿qué campos tiene un mensaje?
\end{lstlisting}

\section{Servicios y parámetros}

\begin{description}
  \item[Servicio:] una pregunta con respuesta (como una llamada telefónica), a diferencia del tópico,
        que es una transmisión continua. El \emph{launch} usa servicios de Gazebo para hacer que la
        cámara de la ventana siga al dron.
  \item[Parámetro:] un ajuste de un nodo, como la altura de búsqueda o la ganancia del controlador.
        Se guardan en archivos \textbf{YAML}:
\end{description}
\begin{lstlisting}
busqueda_incendio:
  ros__parameters:
    altura_busqueda: 30.0     # metros sobre el suelo
    velocidad: 5.0            # m/s
\end{lstlisting}

\section{Archivos launch}

Arrancar diez nodos a mano sería tedioso. Un archivo \textbf{launch} (en Python) los arranca todos
con sus parámetros. El nuestro es \cod{launch/simulacion.launch.py} y recibe \emph{argumentos} para
elegir vehículo, mundo, controlador y misión:
\begin{lstlisting}[style=bash]
ros2 launch cuadrirotor_sim simulacion.launch.py vehiculo:=tarot650 mundo:=bosque_grande controlador:=sta mision:=busqueda
\end{lstlisting}

\section{El puente entre Gazebo y ROS}

Gazebo tiene su propio sistema de mensajes (\emph{gz-transport}). El nodo \cod{ros_gz_bridge}
traduce entre los dos mundos según la lista de \cod{config/puente.yaml}: la cámara, la odometría, la
IMU y el LiDAR van de Gazebo a ROS; las velocidades de los motores y el viento, de ROS a Gazebo.

\section{Paquetes y espacio de trabajo}

El código de ROS se organiza en \textbf{paquetes} (carpetas con un \cod{package.xml}). Los paquetes
viven en un \textbf{espacio de trabajo} (\emph{workspace}); el nuestro es \cod{~/ros2_ws}. Se
compila con \cod{colcon} y, antes de usarlo en cada terminal, se ``activa'' con \cod{source}:

\begin{lstlisting}[style=bash]
cd ~/ros2_ws
colcon build --symlink-install       # compila todos los paquetes de src/
source install/setup.bash            # hace visibles los paquetes en esta terminal
\end{lstlisting}

\begin{aviso}
Si una orden de ROS dice que no encuentra el paquete \cod{cuadrirotor_sim}, casi siempre es porque
olvidaste el \cod{source install/setup.bash} en esa terminal. Para no olvidarlo, agrégalo a
\cod{~/.bashrc}.
\end{aviso}

\begin{ejercicio}
Con la simulación corriendo, abre otra terminal, haz \cod{source} y escribe
\cod{ros2 topic echo /cuadrirotor/odometria --field pose.pose.position}. Verás la posición del dron
cambiar en tiempo real.
\end{ejercicio}
